% ============================================================
%  第四章 系统设计
% ============================================================
\chapter{系统设计}

\section{系统总体架构设计}

\subsection{模块化单体架构}

本系统采用模块化单体架构（Modular Monolith），而非微服务架构。架构选型依据如下：

\begin{enumerate}
  \item \textbf{团队规模}：单人开发场景下，微服务带来的运维复杂度远超其收益；
  \item \textbf{业务耦合}：社团、活动、招募等模块存在较强的业务关联，频繁的跨服务调用会引入额外延迟；
  \item \textbf{演进路径}：模块化单体通过清晰的模块边界，为未来微服务化演进奠定基础。
\end{enumerate}

系统模块依赖关系如图~\ref{fig:module-dep}所示。

\begin{figure}[H]
\centering
\begin{tikzpicture}[
  mod/.style={draw, rectangle, rounded corners, minimum width=2.8cm, minimum height=0.8cm, align=center, font=\small},
  core/.style={draw, rectangle, rounded corners, minimum width=2.8cm, minimum height=0.8cm, align=center, font=\small, fill=blue!15},
  gateway/.style={draw, rectangle, rounded corners, minimum width=10cm, minimum height=0.8cm, align=center, font=\small\bfseries, fill=orange!20},
  arr/.style={-Stealth, thick}
]
\node[gateway] (gw) at (5,6) {community-gateway（入口 / 所有Controller）};

\node[mod] (club)     at (1,4)   {community-club};
\node[mod] (user)     at (3.5,4) {community-user};
\node[mod] (activity) at (6,4)   {community-activity};
\node[mod] (recruit)  at (8.5,4) {community-recruit};
\node[mod] (notice)   at (1,2.5) {community-notice};
\node[mod] (admin)    at (3.5,2.5){community-admin};

\node[core] (core) at (5,1) {community-core（实体 / Repository / JWT / 邮件）};

\draw[arr] (gw) -- (club);
\draw[arr] (gw) -- (user);
\draw[arr] (gw) -- (activity);
\draw[arr] (gw) -- (recruit);
\draw[arr] (gw) -- (notice);
\draw[arr] (gw) -- (admin);

\draw[arr] (club)     -- (core);
\draw[arr] (user)     -- (core);
\draw[arr] (activity) -- (core);
\draw[arr] (recruit)  -- (core);
\draw[arr] (notice)   -- (core);
\draw[arr] (admin)    -- (core);
\end{tikzpicture}
\caption{系统模块依赖关系图}
\label{fig:module-dep}
\end{figure}

\subsection{分层架构设计}

系统后端采用经典四层架构，各层职责如下：

\begin{table}[H]
\centering
\caption{后端分层架构说明}
\begin{tabular}{lll}
\toprule
层次 & 组件 & 职责 \\
\midrule
表现层 & Controller & 接收HTTP请求，参数校验，调用Service，返回响应 \\
业务层 & Service & 业务逻辑处理，事务管理，调用Repository \\
持久层 & Repository & 数据库访问，封装JPA查询 \\
实体层 & Entity & JPA实体映射，数据库表结构定义 \\
\bottomrule
\end{tabular}
\end{table}

\subsection{系统部署架构}

系统整体部署拓扑如图~\ref{fig:deploy}所示。

\begin{figure}[H]
\centering
\begin{tikzpicture}[
  box/.style={draw, rectangle, rounded corners, minimum width=2.5cm, minimum height=0.7cm, align=center, font=\small},
  infra/.style={draw, rectangle, rounded corners, minimum width=2.5cm, minimum height=0.7cm, align=center, font=\small, fill=green!15},
  client/.style={draw, rectangle, rounded corners, minimum width=2.5cm, minimum height=0.7cm, align=center, font=\small, fill=yellow!20},
  arr/.style={-Stealth}
]
\node[client]  (browser) at (0,5)   {浏览器 (Vue 3 SPA)};
\node[box]     (backend) at (4,5)   {Spring Boot\\:8080};
\node[box]     (agent)   at (8,5)   {Python Agent\\:8000};

\node[infra]   (mysql)   at (1,3)   {MySQL 8\\:3306};
\node[infra]   (redis)   at (3.5,3) {Redis\\:6379};
\node[infra]   (mq)      at (6,3)   {RabbitMQ\\:5672};
\node[infra]   (oss)     at (8.5,3) {Aliyun OSS};

\node[infra]   (prom)    at (2,1)   {Prometheus\\:9090};
\node[infra]   (grafana) at (5,1)   {Grafana\\:3000};
\node[infra]   (chroma)  at (8,1)   {Chroma DB};

\draw[arr] (browser) -- node[above,font=\tiny]{HTTP/WS} (backend);
\draw[arr] (backend) -- node[above,font=\tiny]{REST} (agent);
\draw[arr] (backend) -- (mysql);
\draw[arr] (backend) -- (redis);
\draw[arr] (backend) -- (mq);
\draw[arr] (backend) -- (oss);
\draw[arr] (backend) -- node[right,font=\tiny]{metrics} (prom);
\draw[arr] (prom)    -- (grafana);
\draw[arr] (agent)   -- (chroma);
\draw[arr] (agent)   -- (mysql);
\end{tikzpicture}
\caption{系统部署架构图}
\label{fig:deploy}
\end{figure}

\section{数据库设计}

\subsection{ER 图设计}

系统核心实体关系图如图~\ref{fig:er}所示，涵盖用户、社团、活动、招募、通知、财务、资源七大核心实体及其关联关系。

\begin{figure}[H]
\centering
\begin{tikzpicture}[
  entity/.style={draw, rectangle, minimum width=2.2cm, minimum height=0.7cm, align=center, font=\small\bfseries, fill=blue!10},
  rel/.style={draw, diamond, minimum width=1.5cm, minimum height=0.6cm, align=center, font=\tiny, fill=orange!20},
  attr/.style={draw, ellipse, minimum width=1.5cm, minimum height=0.5cm, align=center, font=\tiny},
  line/.style={-}
]
% Entities
\node[entity] (user)     at (0,8)    {t\_user};
\node[entity] (club)     at (4,8)    {t\_club};
\node[entity] (member)   at (4,6)    {t\_club\_member};
\node[entity] (activity) at (8,8)    {t\_activity};
\node[entity] (signup)   at (8,6)    {t\_activity\_signup};
\node[entity] (attend)   at (8,4)    {t\_activity\_attendance};
\node[entity] (batch)    at (0,5)    {t\_recruit\_batch};
\node[entity] (appform)  at (0,3)    {t\_recruit\_application};
\node[entity] (notice)   at (4,4)    {t\_notice};
\node[entity] (finance)  at (4,2)    {t\_club\_finance};
\node[entity] (resource) at (8,2)    {t\_resource\_application};

% Relationships
\draw (user) -- node[above,font=\tiny]{1} (member) node[right,font=\tiny]{N} {};
\draw (club) -- (member);
\draw (club) -- node[above,font=\tiny]{1..N} (activity);
\draw (user) -- node[right,font=\tiny]{1..N} (signup);
\draw (activity) -- (signup);
\draw (signup) -- (attend);
\draw (club) -- node[left,font=\tiny]{1..N} (batch);
\draw (batch) -- (appform);
\draw (user) -- (appform);
\draw (club) -- (notice);
\draw (club) -- (finance);
\draw (club) -- (resource);

% Key attributes (selected)
\node[attr] at (-1.8,8.5) {\underline{id}};
\draw (-1.2,8.5) -- (user);
\node[attr] at (-1.8,7.5) {username};
\draw (-1.2,7.5) -- (user);

\node[attr] at (5.8,8.5) {\underline{id}};
\draw (5.2,8.5) -- (club);
\node[attr] at (5.8,7.5) {name};
\draw (5.2,7.5) -- (club);
\node[attr] at (5.8,7.0) {status};
\draw (5.2,7.0) -- (club);
\end{tikzpicture}
\caption{系统核心ER图}
\label{fig:er}
\end{figure}

\subsection{主要数据表设计}

\subsubsection{用户与角色表}

\begin{table}[H]
\centering
\caption{t\_user 用户表}
\begin{tabular}{llll}
\toprule
字段名 & 类型 & 约束 & 说明 \\
\midrule
id            & BIGINT        & PK, AUTO\_INCREMENT & 用户ID \\
username      & VARCHAR(50)   & UNIQUE, NOT NULL    & 用户名 \\
email         & VARCHAR(100)  & UNIQUE, NOT NULL    & 邮箱 \\
password\_hash & VARCHAR(255)  & NOT NULL            & BCrypt密码哈希 \\
avatar\_url    & VARCHAR(500)  &                     & 头像URL \\
status        & TINYINT       & DEFAULT 1           & 状态(1正常/0禁用) \\
created\_at    & DATETIME      & NOT NULL            & 创建时间 \\
\bottomrule
\end{tabular}
\end{table}

\begin{table}[H]
\centering
\caption{t\_club 社团表}
\begin{tabular}{llll}
\toprule
字段名 & 类型 & 约束 & 说明 \\
\midrule
id            & BIGINT       & PK, AUTO\_INCREMENT & 社团ID \\
name          & VARCHAR(100) & UNIQUE, NOT NULL    & 社团名称 \\
category      & VARCHAR(50)  & NOT NULL            & 社团类型 \\
description   & TEXT         &                     & 社团简介 \\
logo\_url      & VARCHAR(500) &                     & 社团Logo \\
status        & VARCHAR(20)  & NOT NULL            & 状态(PENDING/ACTIVE/FROZEN/DISSOLVED) \\
creator\_id    & BIGINT       & FK(t\_user)         & 创建人ID \\
created\_at    & DATETIME     & NOT NULL            & 创建时间 \\
\bottomrule
\end{tabular}
\end{table}

\subsubsection{活动相关表}

\begin{table}[H]
\centering
\caption{t\_activity 活动表}
\begin{tabular}{llll}
\toprule
字段名 & 类型 & 约束 & 说明 \\
\midrule
id            & BIGINT       & PK, AUTO\_INCREMENT & 活动ID \\
club\_id       & BIGINT       & FK(t\_club)         & 所属社团 \\
title         & VARCHAR(200) & NOT NULL            & 活动标题 \\
location      & VARCHAR(200) &                     & 活动地点 \\
start\_time    & DATETIME     & NOT NULL            & 开始时间 \\
end\_time      & DATETIME     & NOT NULL            & 结束时间 \\
max\_attendees & INT          &                     & 人数上限 \\
checkin\_code  & VARCHAR(20)  & UNIQUE              & 签到码 \\
status        & VARCHAR(20)  & NOT NULL            & 活动状态 \\
\bottomrule
\end{tabular}
\end{table}

\subsection{索引与性能优化设计}

系统针对高频查询场景设计了以下复合索引：

\begin{lstlisting}[language=SQL, caption=关键索引设计]
-- 按社团查询活动列表（高频）
CREATE INDEX idx_activity_club_status
  ON t_activity(club_id, status, start_time DESC);

-- 按用户查询报名记录
CREATE INDEX idx_signup_user_activity
  ON t_activity_signup(user_id, activity_id);

-- 招募申请查询
CREATE INDEX idx_recruit_app_batch_status
  ON t_recruit_application(batch_id, status);

-- 通知未读查询
CREATE INDEX idx_notice_user_read
  ON t_notice_read_status(user_id, is_read);
\end{lstlisting}

\section{安全架构设计}

\subsection{认证流程设计}

登录认证时序图如图~\ref{fig:auth-seq}所示。

\begin{figure}[H]
\centering
\begin{tikzpicture}[
  actor/.style={font=\small},
  line/.style={-Stealth, thin},
  dline/.style={-Stealth, thin, dashed}
]
% Lifelines
\draw[thick] (0,8) -- (0,0);   \node[actor] at (0,8.3) {客户端};
\draw[thick] (3,8) -- (3,0);   \node[actor] at (3,8.3) {AuthController};
\draw[thick] (6,8) -- (6,0);   \node[actor] at (6,8.3) {Redis};
\draw[thick] (9,8) -- (9,0);   \node[actor] at (9,8.3) {MySQL};

% Messages
\draw[line]  (0,7.5) -- node[above,font=\tiny]{POST /auth/login} (3,7.5);
\draw[line]  (3,7.0) -- node[above,font=\tiny]{查询用户} (9,7.0);
\draw[dline] (9,6.7) -- node[above,font=\tiny]{返回用户信息} (3,6.7);
\draw[line]  (3,6.3) -- node[above,font=\tiny]{检查登录失败次数} (6,6.3);
\draw[dline] (6,6.0) -- node[above,font=\tiny]{返回计数} (3,6.0);
\draw[line]  (3,5.5) -- node[above,font=\tiny]{生成AccessToken+RefreshToken} (3,5.2);
\draw[line]  (3,5.0) -- node[above,font=\tiny]{存储RefreshToken} (6,5.0);
\draw[dline] (3,4.5) -- node[above,font=\tiny]{返回AccessToken(Body)+RefreshToken(Cookie)} (0,4.5);

% Token refresh
\draw[line]  (0,3.8) -- node[above,font=\tiny]{POST /auth/refresh (携带Cookie)} (3,3.8);
\draw[line]  (3,3.4) -- node[above,font=\tiny]{验证RefreshToken} (6,3.4);
\draw[dline] (6,3.1) -- node[above,font=\tiny]{Token有效} (3,3.1);
\draw[line]  (3,2.7) -- node[above,font=\tiny]{生成新AccessToken} (3,2.4);
\draw[dline] (3,2.2) -- node[above,font=\tiny]{返回新AccessToken} (0,2.2);
\end{tikzpicture}
\caption{登录认证与Token刷新时序图}
\label{fig:auth-seq}
\end{figure}

\subsection{授权设计}

系统采用基于角色的访问控制（RBAC）模型，接口权限矩阵如表~\ref{tab:rbac}所示。

\begin{table}[H]
\centering
\caption{核心接口权限矩阵}
\label{tab:rbac}
\begin{tabular}{lcccc}
\toprule
接口 & 普通用户 & 社团成员 & 社团管理员 & 系统管理员 \\
\midrule
GET /clubs          & \checkmark & \checkmark & \checkmark & \checkmark \\
POST /clubs         & \checkmark &            &            &            \\
PUT /clubs/\{id\}   &            &            & \checkmark & \checkmark \\
POST /activities    &            &            & \checkmark &            \\
GET /admin/users    &            &            &            & \checkmark \\
POST /admin/approve &            &            &            & \checkmark \\
\bottomrule
\end{tabular}
\end{table}

\subsection{安全防护设计}

\begin{enumerate}
  \item \textbf{登录限流}：基于Redis滑动窗口，同一IP 15分钟内登录失败超过5次，锁定账户；
  \item \textbf{XSS防护}：前端使用Element Plus内置转义，后端通过\texttt{SecurityHeadersFilter}设置\texttt{Content-Security-Policy}响应头；
  \item \textbf{CSRF防护}：采用无状态JWT认证，天然免疫CSRF攻击；
  \item \textbf{SQL注入防护}：全程使用JPA参数化查询，禁止字符串拼接SQL；
  \item \textbf{审计日志}：通过AOP切面拦截敏感操作，记录操作人、时间、IP及操作内容。
\end{enumerate}

\section{接口设计}

\subsection{RESTful API 设计规范}

系统遵循RESTful设计规范，统一响应格式如下：

\begin{lstlisting}[language=Java, caption=统一响应格式]
{
  "code": 200,        // 业务状态码
  "message": "success",
  "data": { ... },    // 响应数据
  "timestamp": 1700000000000
}
\end{lstlisting}

\subsection{核心接口设计}

\begin{table}[H]
\centering
\caption{认证模块核心接口}
\begin{tabular}{llll}
\toprule
方法 & 路径 & 说明 & 权限 \\
\midrule
POST & /api/auth/register  & 用户注册 & 公开 \\
POST & /api/auth/login     & 用户登录 & 公开 \\
POST & /api/auth/refresh   & 刷新Token & 公开 \\
POST & /api/auth/logout    & 登出 & 已认证 \\
POST & /api/auth/forgot-password & 找回密码 & 公开 \\
\bottomrule
\end{tabular}
\end{table}

\begin{table}[H]
\centering
\caption{社团管理核心接口}
\begin{tabular}{llll}
\toprule
方法 & 路径 & 说明 & 权限 \\
\midrule
GET    & /api/clubs          & 社团列表（分页） & 公开 \\
POST   & /api/clubs          & 申请创建社团 & 已认证 \\
GET    & /api/clubs/\{id\}   & 社团详情 & 公开 \\
PUT    & /api/clubs/\{id\}   & 更新社团信息 & 社团管理员 \\
POST   & /api/clubs/\{id\}/members & 添加成员 & 社团管理员 \\
DELETE & /api/clubs/\{id\}/members/\{uid\} & 移除成员 & 社团管理员 \\
\bottomrule
\end{tabular}
\end{table}

\section{AI 服务设计}

\subsection{RAG 问答系统设计}

RAG问答系统整体流程如图~\ref{fig:rag}所示。

\begin{figure}[H]
\centering
\begin{tikzpicture}[
  box/.style={draw, rectangle, rounded corners, minimum width=2.8cm, minimum height=0.7cm, align=center, font=\small},
  arr/.style={-Stealth, thick}
]
\node[box, fill=yellow!20] (query)   at (0,4)   {用户提问};
\node[box]                 (embed)   at (3.5,4)  {Embedding\\向量化};
\node[box, fill=green!15]  (chroma)  at (7,4)    {Chroma\\相似度检索};
\node[box]                 (context) at (7,2)    {相关文档块};
\node[box]                 (prompt)  at (3.5,2)  {Prompt\\拼接};
\node[box, fill=blue!15]   (llm)     at (0,2)    {DeepSeek\\LLM};
\node[box, fill=yellow!20] (answer)  at (0,0)    {生成答案};

\draw[arr] (query)   -- (embed);
\draw[arr] (embed)   -- (chroma);
\draw[arr] (chroma)  -- (context);
\draw[arr] (context) -- (prompt);
\draw[arr] (query)   -- (prompt);
\draw[arr] (prompt)  -- (llm);
\draw[arr] (llm)     -- (answer);
\end{tikzpicture}
\caption{RAG问答系统流程图}
\label{fig:rag}
\end{figure}

\subsection{混合推荐系统设计}

混合推荐系统融合协同过滤与内容推荐两路结果，最终得分计算公式为：

\begin{equation}
  \text{Score}(u, c) = \alpha \cdot \text{Score}_{\text{CF}}(u, c) + (1-\alpha) \cdot \text{Score}_{\text{CB}}(u, c)
\end{equation}

其中，$\alpha$ 为融合权重（默认0.6），$\text{Score}_{\text{CF}}$ 为SVD协同过滤得分，$\text{Score}_{\text{CB}}$ 为语义内容相似度得分。

\section{前端架构设计}

\subsection{页面结构与路由设计}

前端路由采用Vue Router 4，通过导航守卫实现权限控制：

\begin{lstlisting}[language=JavaScript, caption=路由权限守卫]
router.beforeEach((to, from, next) => {
  const userStore = useUserStore()
  if (to.meta.requiresAuth && !userStore.isLoggedIn) {
    next('/login')
  } else if (to.meta.roles && !to.meta.roles.includes(userStore.role)) {
    next('/403')
  } else {
    next()
  }
})
\end{lstlisting}

\subsection{状态管理设计}

用户认证状态通过Pinia Store管理，包含Token、用户信息与权限列表，配合\texttt{pinia-plugin-persistedstate}持久化至\texttt{localStorage}，确保页面刷新后登录状态不丢失。
